\section{Related Work}
\label{sec:related-work}

\alex{Here, the strategy that can works is to introduce and explain the general context to thematic of each section, the existing methods and models and their respective advantages, at end emphasize the literature gap and to discuss how this work fill these gaps.}

\subsection{Synthetic Tabular Health Data}

Synthetic microdata predate deep learning \alex{What is the relationship between microdata with tabular data? I think that lacks an introduction for tabular health-data, because the title of section is about this, but here have no descriptions or concepts for tabular data. Also, the most important is to emphasize that the models mentioned here are related to synthetic tabular data! An objective phrase at the beginning can solve this easily.}. Multiple imputation and probabilistic graphical models generate records from estimated conditional distributions and remain attractive because their assumptions can be inspected \citep{raghunathan2003}. Copulas provide efficient dependence modelling but can struggle when discrete, continuous, multimodal, and rare-event variables interact nonlinearly \citep{wang2021}.

Generative adversarial networks formulate generation as an adversarial game \citep{goodfellow2014}. MedGAN adapted this idea to high-dimensional discrete patient records \citep{choi2017}, and CTGAN addressed mixed-type tables through mode-specific normalization, conditional vectors, and training-by-sampling \citep{xu2019}. Variational autoencoders optimize a likelihood-related lower bound \citep{kingma2014}; TVAE was evaluated alongside CTGAN in the same tabular-synthesis framework \citep{xu2019}. More recently, TabDDPM applied separate diffusion processes to numerical and categorical attributes and outperformed several GAN and VAE baselines \alex{I think that ``extensions'' is better than ``baselines'' here...} on general tabular benchmarks \citep{kotelnikov2023}. The present strengthened analysis uses CTGAN and TVAE to test whether separability-aware selection transfers across generator families; diffusion remains an important future comparator rather than an assumed performance ceiling.

\alex{Here we need to emphasize why your work complement or is better than those ones mentioned here! So, give a little bit more of details (in a objective way)...}

\subsection{Evaluation beyond Aggregate Resemblance}

Health-synthesis studies repeatedly conclude that no single metric establishes usefulness or safety and that evaluation should jointly cover resemblance, downstream utility, and privacy \citep{goncalves2020,hernandez2023,pezoulas2024}. Recent comparisons report dataset- and metric-dependent rankings and a persistent fidelity/utility--privacy trade-off \citep{hernandez2025,nanevski2026}. A review of structured synthetic electronic health records also documents continuing problems with rare events, reproducibility, and privacy evaluation \citep{ghosheh2024}.

Univariate fidelity measures whether synthetic columns resemble their real counterparts, pairwise fidelity evaluates dependence structure, and multivariate distinguishability tests whether a classifier can separate real from synthetic records \alex{It is necessary to add some references.}. None is sufficient for rare-outcome prediction. Aggregate scores can be dominated by majority-class records while the minority conditional distribution is poorly represented \alex{It is necessary to add some references.}. Conversely, a selector can improve classification by producing easy prototypes rather than a faithful minority population \alex{It is necessary to add some references.}.

Train-on-synthetic, Test-on-Real (\tstr{}) evaluation probes whether relations learned from synthetic data transfer to held-out real observations \alex{It is necessary to add some references.}. Real-plus-Synthetic-to-Real evaluation instead addresses augmentation \alex{It is necessary to add some references.}. These regimes answer different questions and must not be pooled. Neither is a privacy test or proof that inferential estimates are unbiased \alex{It is necessary to add some references.}.

\alex{It is need to describe the relation of your work with these existing researches, where your work fill these gaps? This will increase the value of your work.}

\subsection{Class Imbalance, Overlap, and Augmentation}

Accuracy can remain high when minority recall is poor \alex{It is need to put some references!}. Macro-F1 weights both classes equally, while AUPRC is often more informative than AUROC for rare positive outcomes \citep{saito2015}. Conventional methods include random sampling, SMOTE \citep{chawla2002}, cost-sensitive learning, and threshold adjustment \citep{he2009}. These approaches correct frequency or decision cost, but they do not necessarily recover informative class-conditional structure.

A 2026 JMIR \alex{We need to change this for ``A study make a comparison...''. It is not good to make direct mention to the name of journal.} comparison across 13 small health datasets found that sampling with replacement could match more complex generative augmentation in several settings, suggesting that apparent gains may come primarily from increased sample size \citep{huetdastarac2026}. This makes random oversampling a required comparator and sharpens our question: does explicitly controlling generated class geometry add value beyond correcting sample counts?

Class imbalance and class overlap are related but distinct. A small minority class can remain geometrically well separated, whereas a moderately imbalanced outcome can be intrinsically difficult because valid observations occupy the same region of feature space \alex{It is necessary to add some references.}. Adding examples cannot resolve irreducible overlap. Aggressively selecting distant examples may improve an empirical decision boundary while deleting difficult but legitimate cases \alex{It is necessary to add some references.}. The relevant test is therefore whether improved discrimination is accompanied by acceptable calibration and preservation of minority coverage. \alex{This is a good argument to include your discussions: \textit{We need to add some discussions about how this work address this issues!}}

\subsection{Privacy-Utility Tension}

Synthetic data are not anonymous by construction \alex{It is necessary to add some references!}. A high-capacity generator can reproduce or reveal information about its training records \alex{It is necessary to add some references!}. Privacy should be examined through nearest-record, membership-inference, or attribute-inference analyses, or protected through a formal mechanism such as differential privacy \citep{yale2019,dwork2006}. A recent scoping review found substantial heterogeneity in privacy and utility evaluation and warned that common assessments can underestimate risk \citep{kaabachi2025}. LDP-GAN combines local differential privacy with medical-record synthesis \citep{kim2024ldpgan}, while a recent oncology framework evaluates candidate synthetic datasets across several quality dimensions before downstream use \citep{christoforou2025}. These studies support treating utility, fidelity, and privacy as separate outcomes rather than assuming that one implies the others. \alex{Therefore, taking this into account, ... we need to discuss how this work will address this...}

\subsection{Research Gap}

% \alex{The idea of this sections is good, because it can discusses/summarizes the research gap and also aggregate value }

The literature leaves a practical gap at the intersection of rare outcomes, class-overlap diagnosis, transparent post-generation control, and multidataset validation. Generators optimize distribution learning, imbalance methods correct class frequency, and evaluation frameworks usually assess synthetic outputs after generation rather than steering class-conditional geometry. We test whether separability is a distinct, attributable intervention and whether any utility improvement survives checks of calibration, minority structure, fidelity, and privacy risk.

\alex{I think we can add a comparison table for each work mentioned in this section. The table can be divided into different areas where each work addresses existing gaps. Each row of the table can represent a work, and each column can represent a feature that the work covers. You will then mark the feature columns (with a checkmark symbol) for each row or work that includes that feature. The final row will represent this work, and we will mark each column to show that our proposal covers all features.\\

\textbf{NOTE:} I will send you an example of this comparison table to help you build it.\\

Next, we can reference this table and describe their features at end of this section.
}
